\section{Conclusions}
\label{sec:5}
This paper proposes a \myCLO{} scheme to use \myLDPC{} coding together with \myCC{} protocols in a \myUCRS{} to maximize the \mySU{} \mySE{},assuming second-order channel state information is available at the \mySU{} transceiver. The \mySE{} is maximized subject to certain
interference and \myQoS{} constraints using an optimum \myJBPL{} and \myCA{} at the \mySU{} \myTX{}. A \myTHARQ{} type \myII{} is utilized in this system, which aims to improve \mySU{} transmission reliability by combining the signals from subsequent retransmissions using \myEGC{}. In particular, lower and upper bounds on the achievable \mySE{} are comparably tight over a large range of signal-to-noise ratios and the required number of transmissions, obtained over a \myTF{}-selective
fading channel using a constrained optimization approach. If the latter can be done in real time, it might serve as a basis for a more elaborated adaptation beyond the presented basic scheme.
While the current paper discusses a single \mySU{} and a single \myPU{} \myUCRS{} utilizing a \myTHARQ{} type \myII{} with \myCC{}, future research will focus on enhancing the \myTHARQ{} to type \myIII{} utilizing incremental redundancy along with testing the \myCLO{} scheme on a more complex \myUCRS{}.
